% !TEX root = ../main.tex
\corrigesChapitre{8}{Géométrie repérée}
\corrigesSouples


\begin{blocLivre}{cor-8-1}
\begin{corrige}{1}{Lire une équation de droite}
\demarche{Lire les coefficients de la droite, puis tester les points par substitution.}
\begin{questions}
\item Un normal est $\res{(3;-2)}$ ; un directeur est $\res{(2;3)}$. Leur produit scalaire vaut 0.
\item Pour $A$, $3-8+5=0$ : oui. Pour $B$, $0-4+5=1\neq0$ : non.
\end{questions}

\label{figgeo-08-01-cor}
\figureGeo[corrige]{08-01}
\end{corrige}
\end{blocLivre}
\begin{blocLivre}{cor-8-2}
\begin{corrige}{2}{Construire une équation}
\demarche{Un normal et un point permettent d’écrire une équation cartésienne.}
$3(x-2)+4(y+1)=0$, donc $\res{3x+4y-2=0}$. Un directeur est $(-4;3)$. En posant $x=0$, on obtient $y=1/2$ : l’ordonnée à l’origine vaut $\res{1/2}$.
\label{figgeo-08-02-cor}
\figureGeo[corrige]{08-02}
\end{corrige}
\end{blocLivre}
\begin{blocLivre}{cor-8-3}
\begin{corrige}{3}{Écrire un cercle}
\demarche{Comparer les carrés des distances au carré du rayon.}
\begin{questions}
\item $\res{(x-2)^2+(y+1)^2=9}$.
\item Les carrés des distances au centre valent 9, 9 et 5. $A$ et $B$ sont sur le cercle ; $C$ est à l’intérieur strict.
\end{questions}

\label{figgeo-08-03-cor}
\figureGeo[corrige]{08-03}
\end{corrige}
\end{blocLivre}
\begin{blocLivre}{cor-8-4}
\begin{corrige}{4}{Reconnaître une équation de cercle}
\demarche{Compléter les deux carrés avant de lire le centre et le rayon.}
On complète chaque carré :
$x^2-4x=(x-2)^2-4$ et $y^2+6y=(y+3)^2-9$.
L’équation devient $(x-2)^2+(y+3)^2-4-9-12=0$, donc
$\res{(x-2)^2+(y+3)^2=25}$.
Le centre est $\res{\Omega(2;-3)}$ et le rayon $\res{5}$.
Sur la droite $x=2$, il reste $(y+3)^2=25$ :
$y+3=5$ ou $y+3=-5$, donc $y=2$ ou $y=-8$.
Les deux intersections sont $\res{P(2;2)\text{ et }Q(2;-8)}$.
\label{figgeo-08-04-cor}
\figureGeo[corrige]{08-04}
\end{corrige}
\end{blocLivre}
\begin{blocLivre}{cor-8-5}
\begin{corrige}{5}{Deux projections simples}
\demarche{Le pied d’une projection appartient à la droite et le segment de projection lui est perpendiculaire.}
\begin{questions}
\item L’axe des abscisses est horizontal : le projeté garde l’abscisse 3 et prend l’ordonnée 0.
Ainsi $\res{H_1(3;0)}$ et $PH_1=|4-0|=\res{4}$.
\item Un point de la droite $y=x$ s’écrit $H_2(t;t)$.
Le vecteur $\vect{AH_2}=(t-3;t-1)$ doit être orthogonal au directeur $(1;1)$ :
$(t-3)+(t-1)=0$, donc $2t-4=0$ et $t=2$.
Ainsi $\res{H_2(2;2)}$ et $AH_2=\sqrt{(-1)^2+1^2}=\res{\sqrt2}$.
\end{questions}
\label{figgeo-08-05-cor}
\figureGeo[corrige]{08-05}
\end{corrige}
\end{blocLivre}
\begin{blocLivre}{cor-8-6}
\begin{corrige}{6}{Projeter et mesurer une distance}
\demarche{Déplacer le point dans la direction d’un normal, puis imposer l’appartenance à la droite.}
\begin{questions}
\item Le normal est $(2;1)$, donc $H=(4-2\lambda;3-\lambda)$.
Imposer $H\in d$ donne
$2(4-2\lambda)+(3-\lambda)-5=0$, soit $6-5\lambda=0$.
On trouve $\lambda=6/5$.
\item En remplaçant $\lambda$, $\res{H(8/5;9/5)}$.
On vérifie l’appartenance : $2(8/5)+9/5-5=0$.
De plus $\vect{AH}=(-12/5;-6/5)=-(6/5)(2;1)$ est normal à $d$.
Le point satisfait donc les deux conditions d’une projection orthogonale.
\item $AH=\sqrt{(12/5)^2+(6/5)^2}=(6/5)\sqrt5=\res{6/\sqrt5}$.
Cette longueur est la distance du point à la droite.
\end{questions}
\label{figgeo-08-06-cor}
\figureGeo[corrige]{08-06}
\end{corrige}
\end{blocLivre}
\begin{blocLivre}{cor-8-7}
\begin{corrige}{7}{Droites parallèles et perpendiculaires}
\demarche{Comparer les vecteurs normaux pour obtenir les directions des droites.}
\begin{questions}
\item Les normaux de $d$ et $e$ sont proportionnels, donc les droites sont parallèles. Doubler l’équation de $d$ donne un terme constant 2, différent de $-7$ : elles sont distinctes.
\item Le produit des normaux vaut $3\times2-2\times3=0$.
\item On prend pour normal $(2;3)$ : $2(x-1)+3(y-2)=0$, soit $\res{2x+3y-8=0}$.
\end{questions}

\label{figgeo-08-07-cor}
\figureGeo[corrige]{08-07}
\end{corrige}
\end{blocLivre}
\begin{blocLivre}{cor-8-8}
\begin{corrige}{8}{Intersections avec les axes décalés}
\demarche{Substituer l’équation de chaque droite dans celle du cercle.}
\begin{questions}
\item Pour $y=1$, $(x-1)^2=16$, donc $\res{(-3;1)\text{ et }(5;1)}$.
\item Pour $x=1$, $(y+2)^2=25$, donc $\res{(1;3)\text{ et }(1;-7)}$.
\item Le centre est $(1;-2)$, de rayon 5. Les deux points de la seconde question sont les extrémités du diamètre vertical ; les deux autres sont symétriques par rapport à $x=1$.
\end{questions}

\label{figgeo-08-08-cor}
\figureGeo[corrige]{08-08}
\end{corrige}
\end{blocLivre}
\begin{blocLivre}{cor-8-9}
\begin{corrige}{9}{Une famille d'ensembles}
\demarche{Le second membre après complétion des carrés fixe l’existence du cercle.}
\begin{questions}
\item $\res{(x-1)^2+(y+2)^2=5-m}$.
\item Si $m<5$, cercle de centre $(1;-2)$ et de rayon $\sqrt{5-m}$. Si $m=5$, seul le point $(1;-2)$. Si $m>5$, aucun point.
\item Pour $m=1$, le rayon est $\res{2}$.
\end{questions}

\label{figgeo-08-09-cor}
\figureGeo[corrige]{08-09}
\end{corrige}
\end{blocLivre}
\begin{blocLivre}{cor-8-10}
\begin{corrige}{10}{Déterminer une tangente}
\demarche{Le rayon au point de contact est normal à la tangente.}
\begin{questions}
\item $\Omega A^2=3^2+4^2=25$ : $A$ est sur le cercle.
\item Le rayon $(3;4)$ est normal à la tangente : $3(x-4)+4(y-6)=0$, soit $\res{3x+4y-36=0}$.
\item $A$ appartient à la tangente et $\vect{\Omega A}$ lui est normal : ces deux conditions caractérisent le projeté.
\end{questions}

\label{figgeo-08-10-cor}
\figureGeo[corrige]{08-10}
\end{corrige}
\end{blocLivre}
\begin{blocLivre}{cor-8-11}
\begin{corrige}{11}{Une médiatrice}
\demarche{La médiatrice passe par le milieu du segment et lui est perpendiculaire.}
\begin{questions}
\item $\res{I(1;2)}$ et $\vect{AB}=(6;2)$.
\item La médiatrice passe par $I$ et a $\vect{AB}$ pour normal : $6(x-1)+2(y-2)=0$, soit $\res{3x+y-5=0}$.
\item $CA^2=2^2+(-4)^2=20$ et $CB^2=4^2+(-2)^2=20$, donc $CA=CB$. On retrouve aussi $3\times0+5-5=0$.
\end{questions}

\label{figgeo-08-11-cor}
\figureGeo[corrige]{08-11}
\end{corrige}
\end{blocLivre}
\begin{blocLivre}{cor-8-12}
\begin{corrige}{12}{Hauteurs et cercle circonscrit}
\demarche{Les hauteurs se coupent à l’orthocentre ; les médiatrices au centre du cercle circonscrit.}
\begin{questions}
\item La droite $(AB)$ est horizontale, donc la hauteur issue de $C$ est $x=2$.
La hauteur issue de $A$ a pour normal $\vect{BC}=(-4;4)$ : son équation est $-4x+4y=0$, soit $y=x$.
Les deux hauteurs se coupent en $\res{H(2;2)}$.
\item Le pied $K$ est sur $(BC):x+y-6=0$ et sur la hauteur $y=x$.
Ainsi $2x=6$, donc $\res{K(3;3)}$.
\item La médiatrice de $[AB]$ est $x=3$.
Le milieu de $[AC]$ est $(1;2)$ et $\vect{AC}=(2;4)$ est normal à sa médiatrice :
$2(x-1)+4(y-2)=0$, soit $x+2y-5=0$.
Avec $x=3$, on obtient $y=1$ : $\res{\Omega(3;1)}$.
\item Le rayon vérifie $R^2=\Omega A^2=(0-3)^2+(0-1)^2=10$.
L’équation est donc $\res{(x-3)^2+(y-1)^2=10}$.
Les points $H$ et $\Omega$ viennent de deux constructions distinctes : hauteurs pour $H$, médiatrices pour $\Omega$.
\end{questions}
\label{figgeo-08-12-cor}
\figureGeo[corrige]{08-12}
\end{corrige}
\end{blocLivre}
\begin{blocLivre}{cor-8-13}
\begin{corrige}{13}{Une droite coupe un cercle}
\demarche{Résoudre le système droite-cercle par substitution, puis vérifier les points.}
\begin{questions}
\item $(x-2)^2+(3-x)^2=5$, donc $2x^2-10x+8=0$, soit $x^2-5x+4=0$.
\item Les racines sont 1 et 4, donc les points sont $\res{(1;3)\text{ et }(4;0)}$.
\item Pour $\Omega(2;1)$ et le normal $(1;1)$, $\lambda=(2+1-4)/2=-1/2$. Le projeté est $(5/2;3/2)$ ; la distance vaut $\res{1/\sqrt2}<\sqrt5$. La droite est bien sécante.
\end{questions}

\label{figgeo-08-13-cor}
\figureGeo[corrige]{08-13}
\end{corrige}
\end{blocLivre}
\begin{blocLivre}{cor-8-14}
\begin{corrige}{14}{Un cercle et des droites horizontales}
\demarche{Le signe d’un carré égal à une expression de k détermine le nombre d’intersections.}
\begin{questions}
\item Dans l’équation $(x-1)^2+(y-3)^2=9$, on remplace $y$ par $k$ :
$\res{(x-1)^2=9-(k-3)^2}$.
\item Un carré ne peut pas être négatif. Si $9-(k-3)^2>0$, il y a deux racines :
$|k-3|<3\iff0<k<6$.
Si le second membre est nul, il y a une seule racine, pour $k=0$ ou $k=6$.
S’il est négatif, donc pour $k<0$ ou $k>6$, il n’y a aucune intersection.
\item Dans les deux cas de tangence, $x=1$. Les points sont $\res{T_0(1;0)}$ et $\res{T_6(1;6)}$.
\item Pour $k=2$, $(x-1)^2=9-1=8$, donc $x=1\pm\sqrt8=1\pm2\sqrt2$.
On obtient $\res{I(1-2\sqrt2;2)}$ et $\res{J(1+2\sqrt2;2)}$.
\end{questions}
\label{figgeo-08-14-cor}
\figureGeo[corrige]{08-14}
\end{corrige}
\end{blocLivre}
\begin{blocLivre}{cor-8-15}
\begin{corrige}{15}{Une parabole comme lieu géométrique}
\demarche{Écrire l’égalité de deux distances non négatives avant de les élever au carré.}
\begin{questions}
\item Le projeté sur l’axe est $(x;0)$, donc la distance est $|y|$. L’égalité cherchée est $x^2+(y-2)^2=y^2$.
\item En développant, $x^2-4y+4=0$, soit $\res{y=x^2/4+1}$. Les deux distances sont positives ou nulles, donc l’égalité des carrés est équivalente à celle des distances.
\item $5=x^2/4+1$ donne $x^2=16$, donc $\res{(-4;5)\text{ et }(4;5)}$.
\item $f\prime(x)=x/2$, donc la pente en 4 vaut 2 et la tangente est $\res{y=2x-3}$. Pour $M(4;5)$, $MF=\sqrt{16+9}=5$, égal à la distance à l’axe.
\end{questions}

\label{figgeo-08-15-cor}
\figureGeo[corrige]{08-15}
\end{corrige}
\end{blocLivre}
